% --- LaTeX Homework Template - S. Venkatraman ---

% --- Set document class and font size ---

\documentclass[letterpaper, 11pt]{article}

% --- Package imports ---

\usepackage{
  amsmath, amsthm, amssymb, mathtools, dsfont,	  % Math typesetting
  graphicx, wrapfig, subfig, float,                  % Figures and graphics formatting
  listings, color, inconsolata, pythonhighlight,     % Code formatting
  fancyhdr, sectsty, hyperref, enumerate, enumitem } % Headers/footers, section fonts, links, lists

% --- Page layout settings ---

% Set page margins
\usepackage[left=1.35in, right=1.35in, bottom=1in, top=1.1in, headsep=0.2in]{geometry}

% Anchor footnotes to the bottom of the page
\usepackage[bottom]{footmisc}

% Set line spacing
\renewcommand{\baselinestretch}{1.2}

% Set spacing between paragraphs
\setlength{\parskip}{1.5mm}

% Allow multi-line equations to break onto the next page
\allowdisplaybreaks

% Enumerated lists: make numbers flush left, with parentheses around them
\setlist[enumerate]{wide=0pt, leftmargin=21pt, labelwidth=0pt, align=left}
\setenumerate[1]{label={(\arabic*)}}

% --- Page formatting settings ---

% Set link colors for labeled items (blue) and citations (red)
\hypersetup{colorlinks=true, linkcolor=blue, citecolor=red}

% Make reference section title font smaller
\renewcommand{\refname}{\large\bf{References}}

% --- Settings for printing computer code ---

% Define colors for green text (comments), grey text (line numbers),
% and green frame around code
\definecolor{greenText}{rgb}{0.5, 0.7, 0.5}
\definecolor{greyText}{rgb}{0.5, 0.5, 0.5}
\definecolor{codeFrame}{rgb}{0.5, 0.7, 0.5}

% Define code settings
\lstdefinestyle{code} {
  frame=single, rulecolor=\color{codeFrame},            % Include a green frame around the code
  numbers=left,                                         % Include line numbers
  numbersep=8pt,                                        % Add space between line numbers and frame
  numberstyle=\tiny\color{greyText},                    % Line number font size (tiny) and color (grey)
  commentstyle=\color{greenText},                       % Put comments in green text
  basicstyle=\linespread{1.1}\ttfamily\footnotesize,    % Set code line spacing
  keywordstyle=\ttfamily\footnotesize,                  % No special formatting for keywords
  showstringspaces=false,                               % No marks for spaces
  xleftmargin=1.95em,                                   % Align code frame with main text
  framexleftmargin=1.6em,                               % Extend frame left margin to include line numbers
  breaklines=true,                                      % Wrap long lines of code
  postbreak=\mbox{\textcolor{greenText}{$\hookrightarrow$}\space} % Mark wrapped lines with an arrow
}

% Set all code listings to be styled with the above settings
\lstset{style=code}

% --- Math/Statistics commands ---

% Add a reference number to a single line of a multi-line equation
% Usage: "\numberthis\label{labelNameHere}" in an align or gather environment
\newcommand\numberthis{\addtocounter{equation}{1}\tag{\theequation}}

% Shortcut for bold text in math mode, e.g. $\b{X}$
\let\b\mathbf

% Shortcut for bold Greek letters, e.g. $\bg{\beta}$
\let\bg\boldsymbol

% Shortcut for calligraphic script, e.g. %\mc{M}$
\let\mc\mathcal

% \mathscr{(letter here)} is sometimes used to denote vector spaces
\usepackage[mathscr]{euscript}

% Convergence: right arrow with optional text on top
% E.g. $\converge[w]$ for weak convergence
\newcommand{\converge}[1][]{\xrightarrow{#1}}

% Normal distribution: arguments are the mean and variance
% E.g. $\normal{\mu}{\sigma}$
\newcommand{\normal}[2]{\mathcal{N}\left(#1,#2\right)}

% Uniform distribution: arguments are the left and right endpoints
% E.g. $\unif{0}{1}$
\newcommand{\unif}[2]{\text{Uniform}(#1,#2)}

% Independent and identically distributed random variables
% E.g. $ X_1,...,X_n \iid \normal{0}{1}$
\newcommand{\iid}{\stackrel{\smash{\text{iid}}}{\sim}}

% Equality: equals sign with optional text on top
% E.g. $X \equals[d] Y$ for equality in distribution
\newcommand{\equals}[1][]{\stackrel{\smash{#1}}{=}}

% Math mode symbols for common sets and spaces. Example usage: $\R$
\newcommand{\R}{\mathbb{R}}   % Real numbers
\newcommand{\C}{\mathbb{C}}   % Complex numbers
\newcommand{\Q}{\mathbb{Q}}   % Rational numbers
\newcommand{\Z}{\mathbb{Z}}   % Integers
\newcommand{\N}{\mathbb{N}}   % Natural numbers
\newcommand{\F}{\mathcal{F}}  % Calligraphic F for a sigma algebra
\newcommand{\El}{\mathcal{L}} % Calligraphic L, e.g. for L^p spaces

% Math mode symbols for probability
\newcommand{\pr}{\mathbb{P}}    % Probability measure
\newcommand{\E}{\mathbb{E}}     % Expectation, e.g. $\E(X)$
\newcommand{\var}{\text{Var}}   % Variance, e.g. $\var(X)$
\newcommand{\cov}{\text{Cov}}   % Covariance, e.g. $\cov(X,Y)$
\newcommand{\corr}{\text{Corr}} % Correlation, e.g. $\corr(X,Y)$
\newcommand{\B}{\mathcal{B}}    % Borel sigma-algebra

% Other miscellaneous symbols
\newcommand{\tth}{\text{th}}	% Non-italicized 'th', e.g. $n^\tth$
\newcommand{\Oh}{\mathcal{O}}	% Big-O notation, e.g. $\O(n)$
\newcommand{\1}{\mathds{1}}	% Indicator function, e.g. $\1_A$

% Additional commands for math mode
\DeclareMathOperator*{\argmax}{argmax}    % Argmax, e.g. $\argmax_{x\in[0,1]} f(x)$
\DeclareMathOperator*{\argmin}{argmin}    % Argmin, e.g. $\argmin_{x\in[0,1]} f(x)$
\DeclareMathOperator*{\spann}{Span}       % Span, e.g. $\spann\{X_1,...,X_n\}$
\DeclareMathOperator*{\bias}{Bias}        % Bias, e.g. $\bias(\hat\theta)$
\DeclareMathOperator*{\ran}{ran}          % Range of an operator, e.g. $\ran(T) 
\DeclareMathOperator*{\dv}{d\!}           % Non-italicized 'with respect to', e.g. $\int f(x) \dv x$
\DeclareMathOperator*{\diag}{diag}        % Diagonal of a matrix, e.g. $\diag(M)$
\DeclareMathOperator*{\trace}{trace}      % Trace of a matrix, e.g. $\trace(M)$

% Numbered theorem, lemma, etc. settings - e.g., a definition, lemma, and theorem appearing in that 
% order in Section 2 will be numbered Definition 2.1, Lemma 2.2, Theorem 2.3. 
% Example usage: \begin{theorem}[Name of theorem] Theorem statement \end{theorem}
\theoremstyle{definition}
\newtheorem{theorem}{Theorem}[section]
\newtheorem{proposition}[theorem]{Proposition}
\newtheorem{lemma}[theorem]{Lemma}
\newtheorem{corollary}[theorem]{Corollary}
\newtheorem{definition}[theorem]{Definition}
\newtheorem{example}[theorem]{Example}
\newtheorem{remark}[theorem]{Remark}

% Un-numbered theorem, lemma, etc. settings
% Example usage: \begin{lemma*}[Name of lemma] Lemma statement \end{lemma*}
\newtheorem*{theorem*}{Theorem}
\newtheorem*{proposition*}{Proposition}
\newtheorem*{lemma*}{Lemma}
\newtheorem*{corollary*}{Corollary}
\newtheorem*{definition*}{Definition}
\newtheorem*{example*}{Example}
\newtheorem*{remark*}{Remark}
\newtheorem*{claim}{Claim}
\newcommand{\Tau}{\mathcal{T}}
\newcommand{\PSet}{\mathcal{P}}


% --- Left/right header text (to appear on every page) ---

% Include a line underneath the header, no footer line
\pagestyle{fancy}
\renewcommand{\footrulewidth}{0pt}
\renewcommand{\headrulewidth}{0.4pt}

% Left header text: course name/assignment number
\lhead{MAM3000W, 3MS -- Hand-in 2}

% Right header text: your name
\rhead{Alex Cristaudo, CRSALE010}

% --- Document starts here ---

\begin{document}
\subsection*{Question 1}

\begin{enumerate}[label=(\alph*)]
\item True \\Let $(x_n)$ be a convergent sequence in $X$, converging to $x\in X$. Since $X$ is finite, $(x_n)$ is eventually constant (Hand-in 1). That is, there exists some $N\in \N$ such that for all $n \ge N$, $x_n = x$. That means the sequence $(f(x_n))$ is eventually constant, as for all $n\ge N, f(x_n) = f(x)$. So $f(x_n) \to f(x)$ as $n \to \infty$ and thus $f$ is continuous

\item False. \\ Consider metric spaces $(\R, d_E)$ and $(\R, d_D)$, and denote $\R_E$ as the set from the euclidean metric space and $\R_D$ for the discrete space. For $f: \R_D \to \R_E$ defined by $f(x) = x$, clearly $f$ is bijective as $\R _D = \R_E = \R$. Since if $(x_n)$ in $\R_D$ converges to $x \in \R_D$ with respect to $d_D$, the sequence $(x_n)$ must be eventually constant and thus the sequence $(f(x_n))$ must converge to $f(x)$ with respect to $d_E$. So $f$ is continuous. Now we use Proposition 6.32 to prove that $f^{-1}$ is not continuous. The set $(0,1)$ is closed in $\R_D$ (every set is a closed set), but $f((0,1)) = (0,1)$ is not closed in $\R_E$ (for example, sequence $(\frac{1}{n})$ converges to $0 \notin (0,1)$) so $f$ is not closed and thus $f^{-1}$ is not continuous

\item False. \\Let $X = [0,1]$ and $Y = \R$, both equipped with the normal euclidean metric and define $f:[0,1] \to \R$  by $f(x) = x$. Then $d_Y (f(x_1), f(x_2)) = d_Y (x_1, x_2) = d_E(x_1, x_2) = d_X(x_1, x_2) $ for all $x_1, x_2 \in [0,1]$ making $f$ an isometry but is clearly not surjective

\item False. \\Consider $X = Y = [-1,1]$, both equipped with the euclidean metric. Then for $f: [-1,1] \to [-1,1]$ defined by $f(x) = x^2$, we get that for all $x,y \in [-1,1]$, $d_E (f(x), f(y)) = |x^2 - y^2| = |x-y| |x+y| \le |x-y|(|x| + |y|) \le 2 |x-y|$ since $|x|, |y| \le 1$. So $f$ is Libschitz-continuous, but is not injective since $f(-1)=f(1)=1$

\item True. \\Let $\epsilon > 0$. Since $f$ is Lipschitz-continuous, for all $x,y \in X, \ d_Y(f(x), f(y)) \le kd_X (x,y), $ for some constant $k>0$. Let $\delta  = \frac{\epsilon}{k}$. Now let $x,y \in X$ and $d_X(x,y) < \frac{\epsilon}{k}$. Then $d_Y(f(x), f(y)) \le kd_X (x,y) < k \cdot \frac{\epsilon}{k} = \epsilon$ as required and thus $f$ is uniformly continuous

\end{enumerate}
\newpage
\subsection*{Question 2}
Let $X=\{x,y\}$. For $\Tau$ to be a topology, $\Tau$ is a family of subsets of $X$ such that: \\ (T1): $X \in \Tau$ and $\emptyset \in \Tau$ \\ (T2): Any finite intersection of sets in $\Tau$ is still in $\Tau$. \\ (T3): Arbritrary unions of sets in $\Tau$ is still in $\Tau$. \\ Note $\PSet(X) = \{\emptyset, \{x\}, \{y\}, \{x,y\}\}$. All possible families of subsets containing $\emptyset$ and $X$ are: \\
$\Tau _1 = \{\emptyset, \{x,y\}\}$ \\
$\Tau _2 = \{\emptyset, \{x\}, \{x,y\}\}$ \\
$\Tau _3 = \{\emptyset, \{y\}, \{x,y\}\}$ \\ 
$\Tau _4 = \{\emptyset, \{x\}, \{y\}, \{x,y\}\}$ \\ All these sets satisfy (T2) and (T3) so this is a list of all possible topologies of $X$ \\  Notice that $\Tau _1$ and $\Tau _4$ are not homeomorphic to any other topology, since any other topology has a different size and and thus a different finite number of open sets (so any bijection $f$ would map $n$ open sets to $n$ open sets).  Now $\Tau _2$ is homeomorphic to $\Tau _3$ using $f : X \to X$ with $x \mapsto y$ and $y \mapsto x$ . This is then the only homeomorphic topological space

\end{document}